% Lösungen Zufallsexperimente und Pfadregeln (zusammenbau v0.4, Heft)
\einheitenkopf{Zufallsexperimente und Pfadregeln}

\erg{1}{a) einzelnes Ergebnis \quad b) 8 Ergebnisse: (K; K; K), (K; K; Z), (K; Z; K), (K; Z; Z), (Z; K; K), (Z; K; Z), (Z; Z; K), (Z; Z; Z); Laplace: ja, alle gleich wahrscheinlich \quad c) $A \cap B$: ZKK; weder A noch B: KZZ}
\erg{2}{a) $P(L \cup B) = 1 - 0{,}68 = 0{,}32$; $P(L \cap B) = 0{,}25 + 0{,}15 - 0{,}32 = 0{,}08$ \quad b) $P(D \cup A) = 1 - 0{,}6 = 0{,}4$; $P(D \cap A) = 0{,}2 + 0{,}3 - 0{,}4 = 0{,}1$}
\erg{3}{a) ja, Laplace \quad b) $\frac{8}{50} = \frac{4}{25}$ \quad c) Alle Murmeln sind gleich groß und gleich schwer und werden blind gezogen – keine ist bevorzugt, jede hat dieselbe Chance. Weiteres Beispiel: der Wurf einer fairen Münze.}
\erg{4}{a) $\frac{0{,}2n + 2}{n + 2} - 0{,}2 = \frac{1{,}6}{n + 2} > 0$ für jedes n \quad b) $\frac{0{,}3n + 1}{n + 3} - 0{,}3 = \frac{0{,}1}{n + 3} > 0$ für jedes n \quad c) Zähler: die blauen Kugeln von vorher plus die neue blaue; Nenner: alle Kugeln von vorher plus die drei neuen; Laplace: günstige durch alle. \quad d) Zähler: die goldenen von vorher plus die zwei neuen goldenen; Nenner: alle Murmeln von vorher plus die vier neuen; Laplace: günstige durch alle.}
\erg{5}{a) p bleibt gleich \quad b) Äste R $0{,}3$, B $0{,}7$ auf beiden Stufen; $P = 0{,}3 \cdot 0{,}7 = 0{,}21$ \quad c) $0{,}6^3 \cdot 0{,}4^2 \approx 0{,}0346$}

5b)

\baumzwei{R/0{,}3,B/0{,}7}{R/0{,}3,B/0{,}7}{R/0{,}3,B/0{,}7}

\erg{6}{a) $q = 0{,}85^8 \approx 0{,}2725$; $1 - (1 - q)^2 \approx 0{,}471$ \quad b) $q = 0{,}8^6 \approx 0{,}2621$; $1 - (1 - q)^2 \approx 0{,}456$ \quad c) Grün: P $= \frac{2}{36} + \frac{4}{36} = \frac{1}{6}$; Rot: P $= \frac{4}{36} + \frac{1}{36} = \frac{5}{36}$; Pfade (2; 6), (6; 2), (4; 4) bei beiden, der Pasch (4; 4) wiegt bei Grün viermal so viel – die Behauptung ist falsch \quad d) Blau: P $= \frac{4}{36} + \frac{2}{36} = \frac{6}{36}$; Gelb: P $= \frac{2}{36} + \frac{4}{36} = \frac{6}{36}$; Pfade (1; 6), (6; 1), (2; 5), (5; 2) bei beiden – gleich wahrscheinlich, die Behauptung ist falsch}
\erg{7}{a) p ändert sich \quad b) $\frac{10}{25} \cdot \frac{9}{24} = \frac{3}{20}$ \quad c) $\frac{8}{32} \cdot \frac{7}{31} \cdot \frac{6}{30} \cdot \frac{5}{29} \approx 0{,}00195$}
\erg{8}{a) Gleich groß, je $\frac{3}{7}$: Die vier Pfade zum dritten Zug ergeben zusammen $\frac{3}{7}$ – oder Symmetrie: jede Ziehungsposition ist gleich gut. \quad b) Gleich: beide $\frac{4}{12} = \frac{1}{3}$ – ohne Wissen über frühere Züge ist jede Position gleich gut, die Summe aller Pfade zur jeweiligen Position ergibt denselben Wert.}
\erg{9}{a) vorwärts: eine Wahrscheinlichkeit ist gesucht \quad b) $0{,}6 \cdot 0{,}5 + 0{,}4a = 0{,}44$, also $a = 0{,}35$ \quad c) $0{,}7p + 0{,}3(1 - p) = 0{,}5$, also $p = 0{,}5$}
\erg{10}{a) P(A) $= 2p(0{,}6 - p)$, Scheitel bei $p = 0{,}3$; Winkel $= 108^\circ$ \quad b) P(A) $= 2p\left(\frac{2}{3} - p\right)$, Scheitel bei $p = \frac{1}{3}$; Winkel $= 120^\circ$}
